\section{从世界模型到多模态学习}

这场课不是把 Transformer 生硬地套到癌症数据上，而是从一个更通用的问题出发：当现实由多个传感器共同观测、行动又会改变未来时，模型应该怎样形成可用于模拟的内部状态？癌症只是这个问题最困难也最有价值的实例之一。讲义因此先建立 world model（世界模型）与 multimodality（多模态）的操作性定义，再进入具体架构和生物数据。

\subsection{本讲的三段任务与读者契约}

课堂路线由三部分组成：多模态机制、癌症场景、以及仍在推进中的未来系统。这里的顺序很重要：若不先理解“为什么融合”，后面的 H\&E、蛋白和 RNA 只会变成名词列表；若不理解临床问题，架构图又会失去任务约束。读者应始终追踪同一条主线：每种新数据或新模块究竟减少了哪一种不确定性，又是否改变了可执行的实验选择。

\lecturefigure{slide-038-agenda.jpg}{课程路线：多模态机制、癌症语境与未来工作}{00:01:47}

读图时先看三段之间的依赖。第一段提供 translation、disambiguation 和 fusion 的语言；第二段把这些概念映射到肿瘤数据；第三段才讨论更大模型、反事实模拟和可解释性。它不是三个互不相干的专题，而是一条“定义问题—构造数据—训练模拟器—检验边界”的链。

\lecturefigure{slide-040-audience-assumptions.jpg}{听众假设：懂 Transformer 基础，但不预设癌症免疫学背景}{00:02:37}

这张假设页直接决定讲义的粒度。后文会保留 English technical terms，却在第一次出现时解释 H\&E、spatial transcriptomics、imputation、counterfactual 和 calibration。相反，自注意力不会从矩阵乘法的最底层重讲，而是集中解释不同模态进入 $Q/K/V$ 后发生了什么。

\lecturefigure{slide-056-talk-goals.jpg}{讲者希望建立的三层说服目标}{00:04:17}

第一层主张是多模态 Transformer 仍有大量开放设计空间；第二层主张是癌症生物学恰好提供了高维、空间化、跨尺度的数据 substrate；第三层才是 Noetik 是否已经接近药物发现价值。三层证据强度不同：架构机制可以由公式和实验说明，癌症数据的研究价值可以由测量结构说明，而治疗价值必须经过湿实验、动物研究和临床验证。

\teachervoice{00:01:05--00:04:17，老师强调前半段即使脱离癌症也应独立有用；他真正想说服听众的首先是“多模态机器学习仍然很有创造空间”，而不是先接受某个公司的产品叙事。}

\begin{importantbox}{讲义的证据分层}
后文将明确区分三类陈述：\textbf{机制定义}说明模型怎样计算；\textbf{研究结果}说明某个数据集和实验中观察到什么；\textbf{临床主张}需要超出讲座的因果与安全验证。把三者混写，是理解 AI for Science 时最危险的错误之一。
\end{importantbox}

\subsection{世界模型：观测、行动与未来状态}

世界模型并不等于“有一个大的 latent space”。课堂给出的定义强调干预条件：模型接收当前观测，还要能回答采取某个动作后世界将怎样变化。一个只做静态分类的网络可以学习丰富表示，却未必构成可用于规划的世界模型。下面先把观测、动作和未来状态写成条件分布，再解释为什么多个传感器必须在这个分布里共同出现。

\lecturefigure{slide-059-world-model-definition.jpg}{世界模型的操作性定义：由传感器观测与动作共同条件化未来}{00:05:35}

可把图中的关系写成
\[
p_\theta(s_{t+1}\mid o_{\le t}^{(1:M)},a_t),
\]
其中 $s_{t+1}$ 是未来状态，$o_{\le t}^{(1:M)}$ 表示截至时刻 $t$ 来自 $M$ 个模态的观测，$a_t$ 是候选动作，$\theta$ 是模型参数。真正用于决策时，我们比较的不是一个无条件未来，而是不同动作诱导的未来分布。

\begin{importantbox}{World model 的最小判据}
模型至少要回答“在当前证据下，如果采取动作 $a$，未来状态会怎样变化”。只有观测编码，没有行动条件或状态转移，就更适合称为 representation model，而不是完整的世界模型。
\end{importantbox}

\lecturefigure{slide-063-multimodal-perception.jpg}{视觉、音频与文本共同支持行动前的查询}{00:07:00}

圆环路口例子把抽象定义落到直觉：摄像头告诉你车辆位置，麦克风告诉你不可见方向的接近速度，文本可能提供“这个路口高峰期危险”的先验。模型要回答“现在走入路口会发生什么”，必须把不同格式、不同噪声和不同时间尺度的证据合并，而不是简单多数投票。

\lecturefigure{slide-064-modality-fusion.jpg}{多模态学习的核心动作：让不同传感器的信息发生可训练交互}{00:07:17}

黄色箭头表示 fusion（融合）：不同模态不只是各自编码后在末端拼分数，而是允许一个流改变另一个流的表示。若把融合后的隐变量记为 $z$，一个通用写法是
\[
p(y\mid o^{(1)},\ldots,o^{(M)})
=\int p(y\mid z)\,p_\theta(z\mid o^{(1)},\ldots,o^{(M)})\,\mathrm{d}z.
\]
$z$ 不是必须显式存在的单个向量；它可以是一组 token、多个层级的 cross-attention 状态，或由条件归一化调制的整条计算路径。

\begin{knowledgebox}{为什么多模态不是“输入更多”}
输入通道增加并不自动产生多模态推理。关键在于模型是否学会：哪些证据可互相翻译，哪些证据只能共同消除歧义，哪些模态在某个样本中缺失，以及不同传感器的误差如何传播到最终行动。
\end{knowledgebox}

\subsection{本章小结}

本章建立了后续统一语言：世界模型是动作条件下的未来模拟器，多模态学习是让多个观测流形成可用于该模拟的联合状态。癌症任务的难点并非“把图像和表格都喂进去”，而是找到正确的观测层级、融合位置、预测目标与验证闭环。

